\section{Introduction}
\label{sec:intro}

\begin{figure}[tb]
\centering
\includegraphics[width=\textwidth]{signature_mechanism}
\caption{\textbf{Why \ours{} works.}
\textbf{(A)} Training 5 seeds of the same model at $T_\mathrm{max}=100$
and 5 more at $T_\mathrm{max}=150$ reaches nearly identical final
validation accuracy (0.6771 vs.\ 0.6776, $<\!0.001$ apart).
\textbf{(B)} Yet the two endpoints \emph{do not} lie in the same loss
basin: linearly interpolating their weights crosses a $+0.25$~nats
train-loss barrier, and test accuracy dips by $6$~percentage points in
the middle.
\textbf{(C)} On the 13\% of test samples where the 100ep group-consensus
and 150ep group-consensus \emph{disagree}, the mixed ensemble is
\textbf{+10.3~pp} more accurate than either single-length consensus. The
macroscopic $+1.5$~pp ensemble gain is the product of this correction
($10.3\%$) and the disagreement fraction ($13.1\%$).}
\label{fig:hero}
\end{figure}

Deep ensembling \cite{lakshminarayanan2017simple} remains one of the
most effective and widely deployed techniques for boosting accuracy and
uncertainty calibration: train $N$ models with different random seeds,
average their softmax outputs, and pay $N\times$ inference for a
$\sim$1--3~pp accuracy boost. A decade of follow-up work has searched for
stronger diversity sources that surpass seeds alone---snapshot ensembles
via cyclic learning rates~\cite{huang2017snapshot}, Stochastic Weight
Averaging (SWA)~\cite{izmailov2018averaging}, Fast Geometric
Ensembling~\cite{garipov2018loss}, BatchEnsemble~\cite{wen2020batchensemble},
and hyperparameter ensembles~\cite{wenzel2020hyperparameter}. These methods
typically pay for extra diversity with custom schedules, weight averaging
that collapses the predictive distribution, or additional training
infrastructure.

\paragraph{A free diversity axis.}
We identify an additional source of ensemble diversity that standard deep
ensembles completely ignore and that costs essentially nothing to exploit:
the \emph{length} of the cosine schedule itself. Concretely, we train two
groups of seeds with $T_\mathrm{max}=100$ and $T_\mathrm{max}=150$---the
same optimizer, augmentation, seeds, and schedule shape, changing only one
integer---and average softmax outputs across both groups. We call this
trick \ours{}.

This is the minimum possible modification to deep ensembling: no new loss,
no new regularizer, no weight averaging, no custom LR curve. Yet it adds
$+1.5$~pp on EEG emotion classification and $+0.4$--$0.5$~pp on image
benchmarks \emph{beyond} what same-length deep ensembling at matched
compute already gives. Remarkably, a two-run \ours{} ensemble ($N{=}2$
trajectory-length groups) often matches or beats a 10-seed same-length
ensemble: trajectory-length diversity and seed diversity attack different
failure modes.

\paragraph{Why does this work?}
Naively, two runs that end at the same mean accuracy should behave like
two seeds of the same ensemble. We show empirically this intuition is
wrong. A cosine schedule truncated at $T_\mathrm{max}$ makes the
\emph{path} through parameter space depend on $T_\mathrm{max}$, not just
the final accuracy. At epoch 100 of a 150-epoch schedule the LR is still
$\sim\!0.25\cdot \mathrm{lr}_{\mathrm{init}}$---a $250\times$ larger step
than at epoch 100 of a 100-epoch schedule---and the model continues taking
meaningful gradient steps for 50 more epochs before settling. This
low-LR tail is an effective \emph{basin-exploration} phase: the model
wanders into a neighboring local minimum at the same loss but with
different per-sample decision boundaries.

We confirm this directly. On FACED, the 100ep and 150ep trajectories reach
the same mean BAcc ($0.6771$ vs.\ $0.6776$, $<\!0.0005$ apart) yet disagree
on $13\%$ of test samples; on those disagreement samples the mixed
ensemble is $+10.3$~pp more accurate than either group's consensus;
$10.3\% \times 13.1\% = 1.35$~pp predicted gain closely matches the
observed $+1.50$~pp. A weight-space PCA further separates the two groups
into distinct centroids (between/within ratio $3.25$), and a
mode-connectivity probe shows the linear path between a 100ep and a 150ep
checkpoint stays at \emph{chance} for every interpolation coefficient---so
parameter-space averaging fails precisely where prediction-space averaging
succeeds.

\paragraph{Contributions.}
\begin{itemize}[leftmargin=*]
\item \textbf{A new, essentially free ensemble-diversity axis.} \ours{}
adds trajectory length to seed as a second source of ensemble diversity
by changing a single integer in the training recipe---no new loss, no
custom schedule, no weight averaging.
\item \textbf{New FACED EEG emotion SOTA.} Our 10-model \ours{} ensemble
reaches \facedacc{} balanced accuracy, $+5.3$~pp over the EMOD AAAI 2026
baseline ($0.642$) and $+1.1$~pp over a true 10-seed same-length deep
ensemble.
\item \textbf{Cross-modality validation.} The same trick lifts 5 more
benchmarks spanning two vision, one text, and a second EEG dataset
(SEED-V $+3.4\%$, CIFAR-100 $+4.7\%$, CIFAR-10 $+2.4\%$, Fashion-MNIST
$+1.2\%$, MNIST $+0.14\%$, SST-2 $+0.3\%$).
\item \textbf{Head-to-head wins on all four calibration metrics.} Against
SWA, snapshot ensembles, and 10-seed deep ensembles on CIFAR-10 at matched
compute, \ours{} has the best accuracy, NLL, ECE, and Brier score.
\item \textbf{A predictive scaling law.} Across nine benchmarks the
ensemble gain is a power law in $(1 - p_{\mathrm{ind}})$ with exponent
$\alpha=0.85$ (Pearson $r=0.75$, $p=0.019$, $R^2=0.57$)---the trick
delivers the most exactly where ensembling is most needed.
\item \textbf{Mechanistic evidence.} Weight-space PCA (between/within
ratio $3.25$) and linear mode-connectivity probes show the two groups
occupy distinct basins; prediction-space averaging succeeds precisely
because parameter-space averaging across the basins does not.
\item \textbf{Endpoint-ratio is task-robust, not task-tuned.} A direct
sweep $T_B \in \{110, 120, 130, 140, 150, 200\}$ at fixed $T_A=100$
identifies $\rho = T_B/T_A \approx 1.5$ as a universal sweet spot
consistent with the basin-separation analysis.
\end{itemize}
